\section{Related Work}\label{sec:related}

\paragraph{{\it Verification}}  Proving the correctness of distributed
protocols and models is a topic of significant recent
interest~\cite{KA+07,HH+15,SWT18}.  These approaches provide strong
correctness guarantees but require significant investment on the part
of the proof engineer (e.g., by defining suitable inductive
invariants) to complete a verification task~\cite{PM+16,YR+21,MG+19}.
In contrast, our focus in this work is to improve the effectiveness of
falsification techniques - validating the presence of bugs in a
distributed protocol design, rather than their absence.  Athough
\name\ cannot be used to verify the correctness of a model, the burden
we impose on test engineers, providing handler specifications as
PAT{}s, as well as a global safety/liveness property in LTL$_f$, is a
significantly lower hurdle than what is required to verify full
functional correctness of these designs.

\paragraph{{\it Testing}}  Besides P~\cite{DGJ+13,ModP}, which we have
dicussed at length, there have also been other efforts at improving
the capabilities of testing frameworks for distributed systems.
Jepsen~\cite{Jepsen} is randomized testing system that seeks to reveal
bugs when an application is deployed on a weakly-consistent storage
system; Ozkan {\em et al.}~\cite{OM+18} defines a randomized testing
procedure for message-passing distributed systems with guaranteed
lower bounds on the probability of finding a depth-$d$ bug, where $d$
is the minimum length of the sequence of events sufficent to witness
the error.  Morpheus~\cite{YY20} uses partial order sampling and
conflict analysis to control scheduling
decisions. MonkeyDB~\cite{RK+19} uses a demonic scheduling mechanism
to expose safety violations in SQL applications that interact with a
weakly-isolated storage backend.  Clotho~\cite{RK+19} combines static
analysis with a bounded model-checker to generate tests that expose
serializability violations in weakly-consistent database systems.  In
contrast to these efforts, \name\ defines a property-guided synthesis
procedure whose derived controller is specialized based on the
constraints defined by handler specifications. In this sense, our
approach can be seen as a form of property-based testing
(PBT)~\cite{CH00,HC+24} applied to open reactive systems.  Broadly
related to our approach is Mocket~\cite{WD+23}, a PBT-style testing
framework that uses the state space graph extracted from
model-checking TLA+ specifications~\cite{Lam02} to force executions to
follow specific paths in the graph.  Unlike \name, Mocket requires
manual instrumentation of implementations to align actions defined in
the specification with the corresponding code in the implementation,
and relies on the TLC model-checker to produce the state space graph.
On the other hand, \name\ uses a compositional refinement type system
to drive synthesis, and requires no instrumentation or \emph{a priori}
enumeration of the state space to synthesize its controllers.

\paragraph{{\it Specifications}}
TLA+~\cite{Lam02} is a specification language based on LTL for
modeling distributed systems; the correctness of these specifications
are verified using the TLC model checker, and has had notable
real-world impact~\cite{NR+15}.  While \name{}'s use of LTL$_f$
specifications in \PAT{}s is a point of commonality with TLA+, the
integration of these specifications within a refinement type system,
and their role in driving a component-based synthesis procedure
differentiates \name{}'s motivation and design from TLA+ and TLC in
obvious ways.  Type and effect systems that target \emph{temporal}
properties on the sequences of effects a program may \emph{produce} is
a well-studied subject. For example, \cite{SS+04} presents a type and
effect system for reasoning about the shape of histories (i.e., finite
traces) of events embedded in a program.  Koskinen and
Terauchi\cite{KT+14} present a type and effect system that
additionally supports verification properties of infinite traces,
specified as B\"{u}chi automata.  More recently,
\cite{Temporal-Verification} consider how to support richer control
flow structures, e.g., delimited continuations, in such an effect
system.  Closest to our work, is the proposal of Zhe {\it et
  al.}~\cite{ZYDJ24} on \emph{Hoare Automata Types} (HAT), an
integration of symbolic finite automata within a refinement type
system, suitable for reasoning about stateful sequential programs
structured as a functional core interacting with opaque effectful
libraries.  \PAT{}s, while sharing commonality with HATs, extends them
in important ways, most notably with its incorporation of prophency
automata to enable their use in a distributed setting, where traces
represent both constraints on the history of consumed messages as well
as requirements of future messages that have yet to be handled.

\section{Conclusion}\label{sec:conc}

This paper examines a property-guided testing framework for open
reactive distributed system models.  Our key innovation is the use of
prophecy automata types (\PAT{}s) that enable the specification of
message handlers in terms of a trace history and future.  A
component-based synthesis procedure leverages \PAT{}s to output
bespoke test controllers specialized to generating execution traces
that violate a given safety or liveness property.  Experimental
results on a wide range of benchmarks, including real-world models
used in production, show that \name\ is significantly more effective
in uncovering design bugs than the existing state-of-the-art.
